\chapter{Chapter 13: What Simulations Can and Cannot Tell Us}\label{chapter-13-what-simulations-can-and-cannot-tell-us}

\emph{In which we reckon honestly with the limits of our own work.}

\begin{center}\rule{0.5\linewidth}{0.5pt}\end{center}

\section{The Seduction of Emergence}\label{the-seduction-of-emergence}

When a simulation produces a pattern that resembles reality, the temptation is immediate: this \emph{explains} reality. Our engine produces cooperation from thermodynamics, and cooperation exists in nature, therefore thermodynamics explains cooperation in nature.

This reasoning is wrong, and we must say so clearly.

Our simulation demonstrates \emph{sufficiency}, not \emph{necessity}. Thermodynamic enforcement with FEP gradient descent is \emph{sufficient} to produce cooperation. It is not the \emph{only} mechanism that can produce cooperation. Game theory produces cooperation. Kin selection produces cooperation. Reciprocal altruism produces cooperation. Cultural evolution produces cooperation. Our engine adds thermodynamic enforcement to this list --- it does not replace the others.

\section{What We Can Claim}\label{what-we-can-claim}

\textbf{Claim 1: Cooperation can emerge from physics alone.} No strategy, no learning, no communication, no intelligence required. This is a constructive proof. We built a system where cooperation emerges from energy constraints, and it works across 45 findings. The claim is narrow: it is \emph{possible} for cooperation to have purely physical origins.

\textbf{Claim 2: The specific integration metric doesn't matter.} IIT $\Phi$, Shannon entropy, and constant scalars produce identical cooperation. Only zero integration differs. This is a universality claim about the model --- the pattern depends on the cost structure, not the internal metric.

\textbf{Claim 3: The critical parameters are social range and mutual benefit rate.} These are CRITICAL (d \textgreater{} 1.0) in sensitivity analysis. Other parameters (ambient regen, maintenance cost) matter less. This identifies the engine's leverage points.

\textbf{Claim 4: A phase transition separates cooperative and collapsed regimes.} This is a mathematical observation about the model's dynamics. The transition is sharp, bistable above threshold, and governed by maintenance pressure.

\textbf{Claim 5: Voluntary cooperation can decline; involuntary cooperation cannot.} Adding cooperation willingness reproduces Putnam's four trends. Removing it makes cooperation immune to decline.

\section{What We Cannot Claim}\label{what-we-cannot-claim}

\textbf{We cannot claim our agents are conscious.} $\Phi$ is a coupling parameter, not a measurement of subjective experience. Chapter 3 classifies it as INSPIRED at best.

\textbf{We cannot claim these results apply to human societies.} Our agents have no culture, no institutions, no language, no history, no identity. The structural parallels to Putnam, Klinenberg, Merton, and the resource curse are \emph{suggestive}, not \emph{explanatory}. The correct framing is: ``a thermodynamic model produces a pattern that structurally resembles X'' --- not ``thermodynamics explains X.''

\textbf{We cannot claim the FEP gradient implements the Free Energy Principle.} Friston's FEP is a variational principle derived from Bayesian inference under a generative model. Our gradient is a handcrafted weighted sum of four heuristic components. It captures the \emph{spirit} of FEP (agents move to reduce surprise) but not the \emph{mathematics} (there is no generative model, no variational bound, no Bayesian updating).

\textbf{We cannot claim our thermodynamics is complete.} We track U, T, S, and F, but our agents are not thermodynamic systems in the physics sense. They have no heat bath, no thermal equilibrium, no partition function. The second law is enforced by code (entropy monotonically increases), not derived from statistical mechanics.

\textbf{We cannot claim the negative results are powered to detect small effects.} Twenty seeds at each condition provides adequate power for large effects (d \textgreater{} 0.8) but may miss small-to-medium effects. The null results (information, memory, communication, dimensionality) could reflect genuinely zero effects or insufficient statistical power. We report effect sizes throughout and let the reader judge.

\section{The Value of Constructive Proofs}\label{the-value-of-constructive-proofs}

Despite these limitations, the simulation has scientific value as a \emph{constructive proof} --- a demonstration that a specific set of assumptions produces a specific set of outcomes.

The cooperation emergence finding (Chapter 5) is a constructive proof that energy scarcity + FEP gradient + harmony resonance $\to$ spatial clustering + mutual survival. The ablation (Finding 2) decomposes this into necessary components. The null model (Finding 32) establishes that specific mechanisms --- not just agent movement --- are required. The falsification attempt (Finding 38) demonstrates robustness across parameter space.

These are not explanations of nature. They are existence proofs for a class of dynamical systems. The value is in the existence, not the explanation.

\section{The Honest Middle}\label{the-honest-middle}

There is a position between ``this simulation explains reality'' (overclaim) and ``this simulation is just a toy'' (underclaim). The honest middle is:

\emph{This simulation demonstrates that a specific set of physical rules produces emergent patterns that structurally resemble observed sociological phenomena. The resemblance is not evidence of shared mechanism --- human cooperation involves cognition, culture, and choice that our agents lack. But the resemblance is also not coincidence --- both systems face the same structural challenge (finite resources, spatial constraints, mutual benefit from proximity) and both systems produce the same structural response (clustering, cooperation, phase transitions).}

The structural challenge is real. The structural response is real. The mapping between the two is suggestive, not causal.

We leave causal claims to those with experimental access to human societies. We offer, instead, a thermodynamic framework that generates the right \emph{shapes} --- and an honest accounting of everything that framework cannot do.

\begin{center}\rule{0.5\linewidth}{0.5pt}\end{center}

\emph{Next: Chapter 14 --- The Equal Opportunity Paradox and the Freedom to Bowl Alone}
